\documentclass[10pt,a4paper]{amsart}
\usepackage[margin=19mm]{geometry}
\usepackage{amsmath,amssymb,mathtools}
\usepackage[scr=rsfs,cal=euler]{mathalfa}
\usepackage{xfrac}
\usepackage[unicode,hidelinks,bookmarks=false]{hyperref}
\newcommand{\ie}{i.e.\ }
\newcommand{\cf}{cf.\ }
\newcommand{\resp}{respectively\ }
\newcommand{\Subsection}{\subsection}
\providecommand{\nobreakdash}{\nobreak\hbox{-}}
\setlength{\emergencystretch}{2em}
\multlinegap0pt
\usepackage{bm}
\usepackage{bbm}
\usepackage{dsfont}
\usepackage{xspace}
\usepackage{cancel}
\usepackage{mathtools}
\usepackage[shortlabels]{enumitem}
\usepackage{amsthm}
\makeatletter
\@ifundefined{theorem}{\newtheorem{theorem}{Theorem}}{}
\makeatother
\newcommand{\myabs}[1]{ | #1 | }
\newcommand{\myinnerproduct}[2]{ \langle #1, #2 \rangle }
\title[On a comprehensive review of a proof of L\"owner's theorem]{On a comprehensive review of a proof of L\"owner's theorem}
\author{Curt Healey}
\date{Source: arXiv:2510.00121v2 (CC BY 4.0)}
\begin{document}
\maketitle

\noindent\emph{Source and licence.} On a comprehensive review of a proof of L\"owner's theorem. Version arXiv:2510.00121v2 (CC BY 4.0), distributed under the Creative Commons Attribution 4.0 International licence, as recorded in the arXiv licence metadata for this exact version and shown on the abstract page https://arxiv.org/abs/2510.00121v2. Full rights notice: licenses/arxiv-2510.00121-CC-BY-4.0.txt.\par\smallskip

\noindent\textit{Editorial scope.} Counted unit: the Theorem whose source label is \texttt{first-divided-difference-positive-implies-matrix-monotone} in the article (source0.tex lines 661--664, source label \texttt{first-divided-difference-positive-implies-matrix-monotone}), delivered together with the article's own complete original proof of it (source0.tex lines 665--721, one \texttt{proof} environment of 57 lines). No mathematics is written, repaired or re-derived: statement and proof are the article's own text line for line. Required context retained uncounted: chain-rule-for-matrices (lines 643--658). Every definition, display and label of those lines is retained unchanged. Omitted from this package and not redistributed: 1--642, 659--660, 722--1078. Nothing inside a retained excerpt is omitted or altered: every retained body is delivered exactly as the article's lines are written, so no omission receipt of any kind accompanies this package. Citations: the retained excerpts carry no citation command at all, so no bibliography entry of the article is transcribed and no reference is redistributed. Reference adaptation: none was needed and none was made; every reference inside the delivered unit resolves inside it, so no unresolved reference or citation is produced. Printed slips kept literal: no printed word, symbol, hypothesis or number of the counted statement or of its proof was altered. Layout and class: the article's own class is \texttt{amsart} with packages including amssymb, amscd, latexsym, amsxtra, dsfont, enumerate, changes, amsfonts, bbold, eucal, calrsfs, enumitem, fge, mathtools; the wrapper is the lane's pinned \texttt{amsart} editorial wrapper at 10pt on A4, which restates the theorem-like environments the retained text needs and adds the article's own macro definitions that the retained text needs (\texttt{\textbackslash myabs}, \texttt{\textbackslash myinnerproduct}), with extra wrapper packages (bm, bbm, dsfont, xspace, cancel, mathtools, enumitem). The delivered numbering is the unit's own. Assets: no retained span contains a figure environment, a graphics-inclusion command, a PDF-page inclusion, a table or a TikZ picture; no image file is copied into this package, the package declares no asset dependency and no image file is redistributed with it. Licence: version arXiv:2510.00121v2 of this article is distributed under the Creative Commons Attribution 4.0 International licence, as recorded in the arXiv licence metadata for this exact version and shown on the abstract page https://arxiv.org/abs/2510.00121v2. A full rights notice is delivered at licenses/arxiv-2510.00121-CC-BY-4.0.txt. Characters: every retained excerpt is pure ASCII, so the delivered engine, XeLaTeX with its default Latin Modern text font, reports no missing character.
	\begin{theorem}[Chain rule]\label{chain-rule-for-matrices}
		Let $f :(a, b) \rightarrow \mathbb{R}$ and $\gamma: (c,d) \rightarrow H_n(a, b)$ be functions on $(a, b), (c,d) \subseteq \mathbb{R}$ such that $\textnormal{sp}(\gamma(t)) \subseteq (a,b)$ for all $t \in (c, d)$. Let $\gamma(t) = U(t) D(t) U(t)^{*}$, where $U$ is a unitary matrix, and $D=\textnormal{diag}(\lambda_1(t),  \ldots, \lambda_n(t))$ for all $t \in (c,d)$. Then, the following hold
		\begin{enumerate}[label=\textnormal{(\Roman*)}]
			\item If $f$ and $\gamma$ are continuously differentiable, then $f \circ \gamma$ is continuously differentiable with derivative 
			\[ \frac{d}{dt} f(\gamma(t)) = U(t) \bigg( [\Delta f(\lambda_i(t), \lambda_j(t))]_{i, j = 1}^{n} \circ [U(t)^{*} \gamma'(t)U(t)] \bigg) U(t)^{*}  \]
			for all $t \in (c, d)$.
			\item  If $f$ and $\gamma$ are twice continuously differentiable, then $f \circ \gamma$ is twice continuously differentiable. Let $c_k(t)$ denote the $k$'th column of $U(t)^{*}\gamma'(t)U(t)$, then the second derivative of $f \circ \gamma$ is
			\begin{equation*}
				\begin{split}
					\frac{d^{2}}{dt^{2}} f(\gamma(t)) = U(t) &\bigg( 2 \sum_{k=1}^{N} [\Delta^{2}f(\lambda_i(t), \lambda_j(t), \lambda_k(t))]^{n}_{i,j=1} \circ (c_k(t) c_k^{*}(t))   \\
					&+[\Delta f(\lambda_i(t), \lambda_j(t))]_{i, j = 1}^{n} \circ (U(t)^{*} \gamma''(t) U(t)) \bigg)U(t)^{*}
				\end{split}
			\end{equation*}
			for all $t \in (c, d)$.
		\end{enumerate}
	\end{theorem}

	\begin{theorem}\label{first-divided-difference-positive-implies-matrix-monotone}
		Let $f:(a,b) \rightarrow \mathbb{R}$ be a continuously differentiable function on $(a, b ) \subseteq \mathbb{R}$. Then $f$ is operator monotone of order $n$ iff 
		\[  [\Delta f(t_i, t_j)]_{i,j=1}^{n} \geq 0, \quad \forall \, t_1, \ldots,  t_n \in (a, b).\]
	\end{theorem}

	\begin{proof}
		Suppose the matrix of first divided differences is positive semi-definite. Let $A, B \in H_n(a, b)$ and $A \leq B$. Let $\epsilon_1 > 0$ such that $\text{sp}((1-t)A + tB) \subseteq (a, b)$ for $t \in (-\epsilon_1, 1 + \epsilon_1)$. Consider $\gamma(t) = (1-t)A + tB$ for $t \in (-\epsilon_1, 1 + \epsilon_1)$. Furthermore, $\gamma'(t) = B-A$ for $t \in (-\epsilon_1, 1 + \epsilon_1)$, so by (I) of Theorem \ref{chain-rule-for-matrices}, we have 
		\begin{align*}
			\begin{split}
				&f(B)-f(A) = \int_{0}^{1} \frac{d}{dt} f(\gamma(t)) \, \rm{d}t \\
				&= \int_{0}^{1} U(t) \bigg( [\Delta f(\lambda_i(t), \lambda_j(t))]_{i,j=1}^{n} \circ (U(t)^{*} (B-A) U(t)) \bigg)U(t)^{*} \, dt.
			\end{split}
		\end{align*}
		Since  $[\Delta f(\lambda_i(t), \lambda_j(t))]_{i,j=1}^{n}  \geq 0$ and $U(t)^{*} (B-A) U(t) \geq 0$, the Schur product is positive semi-definite as well. This implies that the integral is the limit of sums of positive semi-definite matrices, so $f(A) \leq f(B)$.
		\par 
		Suppose $f$ is operator monotone of order $n$. Let $x = (x_1, \ldots, x_n) \in \mathbb{C}^{n}$ and let $D = \text{diag}(t_1, \ldots, t_n)$ where $\lbrace t_1, \ldots, t_n \rbrace \subseteq (a, b)$. Let $\epsilon_2 > 0$ such that $A = D - \epsilon_2 (xx^{*})$ and  $B  = D + \epsilon_2 (xx^{*})$ are elements of $H_n(a,b)$. Let $\gamma(t) = (1-t) A + tB$ for $t \in (0,1)$. Let us now apply (I) of Theorem \ref{chain-rule-for-matrices} for $f'(\gamma(1/2))$.  Since  $\gamma(1/2)$ is a diagonal matrix, it follows that $U(1/2) = \mathbb{I}_n$. Furthermore,  $\gamma^{'}(1/2)= 2 \epsilon_2 (x x^{*})$. These two observations imply that
		\begin{equation*}\label{horn-th-6.6.36-eq1}
			(f \circ \gamma)'(1/2)  = \frac{d}{dt} f(\gamma(t)) \vert_{t=1/2} = [\Delta f(\lambda_i, \lambda_j)]_{i,j=1}^{n} \circ 2 \epsilon_2(x x^{*}).
		\end{equation*}
		
		
		%		 		\Delta f(\lambda_2, \lambda_1) x_2 x_1 & \Delta f(\lambda_2, \lambda_2) \myabs{x_2}^{2} & \ldots & \Delta f(\lambda_2, \lambda_n) x_2 x_n \\
		
		Moreover, if $t \geq s$ for $t, s \in (0,1)$, then 
		\[  \gamma(t) - \gamma(s) =  (1-t)A + tB - (1-s)A + sB   = (t-s)(B-A).  \]
		However, $B-A = 2 \epsilon_2 xx^{*}$ is positive semi-definite since $\myinnerproduct{(xx^{*})y}{y} = y^{*}xx^{*}y = (y^{*}x)(y^{*}x)^{*} = \myabs{y^{*}x}^{2}$. Thus,  $\gamma(t) \geq \gamma(s)$ and since $f$ is operator monotone of order $n$, then $f(\gamma(t)) \geq f(\gamma(s))$. This allows us to conclude that
		\[ \frac{d}{dt} f(\gamma(t)) \vert_{t=1/2} = \lim_{h \rightarrow 0+} \frac{f(\gamma(1/2 + h)) - f(\gamma(1/2))}{h} \]
		is positive semi-definite. Consider $y = (1, \ldots, 1) \in \mathbb{R}^{n}$. Then 
		\begin{equation}\label{first-divided-difference-positive-implies-matrix-monotone-eq1}
			0 \leq \bigg \langle \frac{(f \circ \gamma)'(1/2)}{2 \epsilon_2} y, y  \bigg \rangle . 
		\end{equation}
		Computing $y^{*}([\Delta f(\lambda_i, \lambda_j)]_{i,j=1}^{n} \circ (x x^{*}))y$, we get
		\begin{align*}
			\begin{pmatrix}
				1, \ldots, 1
			\end{pmatrix} 
			&\begin{bmatrix}
				\Delta f(\lambda_1, \lambda_1) x_1 x^{*}_1 & \Delta f(\lambda_1, \lambda_2) x_1 x^{*}_2 & \ldots & \Delta f(\lambda_1, \lambda_n) x_1 x^{*}_n \\
				\vdots & \vdots & \vdots & \vdots \\
				\Delta f(\lambda_n, \lambda_1) x_n x^{*}_1 & \Delta f(\lambda_n, \lambda_2) x_n x^{*}_2 & \ldots & \Delta f(\lambda_n, \lambda_n) x_n x^{*}_n
			\end{bmatrix}
			\begin{pmatrix}
				1 \\
				\vdots \\
				1
			\end{pmatrix} 
			\\
			&= 		
			\begin{pmatrix}
				1, \ldots, 1
			\end{pmatrix} 
			\begin{pmatrix}
				x_1 \sum_{i=1}^{n} \Delta f(\lambda_1, \lambda_i) x^{*}_i \\
				\vdots \\
				x_n \sum_{i=1}^{n} \Delta f(\lambda_n, \lambda_i) x^{*}_i
			\end{pmatrix} 
			= \sum_{j=1}^{n} x^{*}_j \sum_{i=1}^{n} \Delta f(\lambda_j, \lambda_i) x_i
			\\
			&= \myinnerproduct{[\Delta f(\lambda_i, \lambda_j)]_{i,j=1}^{n} x}{x}.
		\end{align*}		
		Therefore, by combining the above and (\ref{first-divided-difference-positive-implies-matrix-monotone-eq1}), we conclude that \[[\Delta f(\lambda_i, \lambda_j)]_{i,j=1}^{n} \circ (x x^{*}) \] is positive semi-definite.
	\end{proof}

\end{document}
